\chapter{Doctoral Entrance Problems}
\label{ch:problems}

The problems below are the algebra of Algerian doctoral entrance examinations,
reproduced in French as they were set. They are grouped by university and, inside
each university, ordered from the oldest paper to the most recent. Every problem
carries the university, the year and the number of the épreuve it came from.

Each statement is given exactly as it was examined: no hint has been added and no
step has been removed. Where the original scan is damaged, the source line says
so rather than silently inventing the missing symbols. Full solutions are in
\cref{ch:solutions}.
\newpage

%% ===================================================================
\section{Centre Universitaire Si El Haouès --- Barika}
%% ===================================================================

\begin{problem}{(Barika) — Épreuve N°02}{alg:barika-2025-e2-p1}
Let $(G, 1_G, \leq)$ be an ordered group with the identity element $1_G$.
\begin{enumerate}
\item Prove that if $G$ is a $\vee$-semilattice, then for all $x, y, z \in G$ it
      holds that
      $$
      x(y \vee z) = xy \vee xz \quad \text{and} \quad (y \vee z)x = yx \vee zx.
      $$
\item Prove that the following statements are equivalent:
      \begin{enumerate}
      \item $G$ is a $\vee$-semilattice;
      \item $G$ is a lattice (indication: show that $x \wedge y = x(x \vee y)^{-1}y$
            for any $x, y \in G$);
      \item $\sup\{x, 1_G\}$ exists for any $x \in G$.
      \end{enumerate}
\item Suppose that $G$ is a lattice and $H$ is a subgroup of $G$. Prove that $H$ is
      a sublattice of $G$ if and only if for all $x \in H$, $\sup\{x, 1_G\} \in H$.
\end{enumerate}

\source{\emph{\small Centre Universitaire Si El Haouès - Barika, 2025, Épreuve~n\textsuperscript{o}~02 — barème : 7 points.}}
\end{problem}

\newpage

\begin{problem}{(Barika) — Épreuve N°02}{alg:barika-2025-e2-p2}
Let $G$ be a group with identity element $1_G$ and $P$ a subset of $G$ satisfying
the following three conditions:

(i) $P \cap P^{-1} = \{1_G\}$; (ii) $P^2 = P$, where $P^2 = \{xy \mid x, y \in P\}$;
(iii) $\forall x \in G$, $xPx^{-1} = P$.

Consider the binary relation $\leq$ defined on $G$ by
$x \leq y \Leftrightarrow yx^{-1} \in P$.
\begin{enumerate}
\item Prove that $(G, \leq)$ is an ordered group.
\item Deduce that $P$ is the positive cone relative to this compatible order on $G$.
\item Prove that $\leq$ is total if and only if $P \cup P^{-1} = G$.
\item \textbf{Application:} consider the additive abelian group
      $\mathbb{R} \times \mathbb{R}$ and the subset $P$ defined by
      $$
      (x, y) \in P \Leftrightarrow (x = y = 0) \text{ or } (x \geq 0 \text{ and } y > 0).
      $$
      \begin{enumerate}
      \item Prove that $P$ is the positive cone of the compatible order $\leq$ on
            $\mathbb{R} \times \mathbb{R}$.
      \item Show that this order is described by
            $$
            (x, y) \leq (x', y') \Leftrightarrow
            \begin{cases} x = x' \text{ and } y = y'; \\
            \text{or } x \leq x' \text{ and } y < y'. \end{cases}
            $$
      \end{enumerate}
\end{enumerate}

\source{\emph{\small Centre Universitaire Si El Haouès - Barika, 2025, Épreuve~n\textsuperscript{o}~02 — barème : 7 points.}}
\end{problem}

\newpage

\begin{problem}{(Barika) — Épreuve N°02}{alg:barika-2025-e2-p3}
Let $G$ be a group and $H$ be a fuzzy subgroup of $G$. $H$ is called a normal
fuzzy subgroup if for all $x, y \in G$, $\mu_H(xy) = \mu_H(yx)$, where $\mu_H(x)$
is the degree of membership of $x$ in the fuzzy set $H$.
\begin{enumerate}
\item Prove that $H$ is a fuzzy normal subgroup of $G$ if and only if $\mu_H$ is
      constant in each conjugate class of $G$.

      \emph{Note that} $\mu_H$ constant in each conjugate class of $G$ means
      $\mu_H(yxy^{-1}) = \mu_H(x)$ for all $x, y \in G$.
\item Prove that if $H$ is a fuzzy normal subgroup, then
      $$
      \mu_H[x, y] \geq \mu_H(x) \quad \text{for all } x, y \in G,
      $$
      where $[x, y] = x^{-1}y^{-1}xy$ is the commutator of $x$ and $y$.
\item Let $\alpha \in [0, 1]$ be such that $\alpha \leq \mu_H(e)$, where $e$ denotes
      the identity of $G$. Prove that if $H$ is a fuzzy normal subgroup of $G$, then
      $H_\alpha = \{x \in G \mid \mu_H(x) \geq \alpha\}$ is a normal subgroup of $G$.
\item Let $\varphi : G \to G'$ be a group morphism and $B$ a fuzzy set of $G'$.
      Consider the fuzzy set $\varphi^{-1}(B)$ of $G$, called the preimage of $B$,
      defined by
      $$
      \forall x \in G,\; \mu_{\varphi^{-1}(B)}(x) = \mu_B(\varphi(x)).
      $$
      Prove that $\varphi^{-1}(B)$ is a fuzzy normal subgroup of $G$.
\end{enumerate}

\source{\emph{\small Centre Universitaire Si El Haouès - Barika, 2025, Épreuve~n\textsuperscript{o}~02 — barème : 6 points. Dans l'énoncé original la question 3 écrit « $x \in L$ », probablement une coquille pour « $x \in G$ ».}}
\end{problem}

\newpage

%% ===================================================================
\section{École Normale Supérieure de Kouba (ENS)}
%% ===================================================================

\begin{problem}{(ENS Kouba) — Épreuve N°03}{alg:ens-kouba-2023-e3-p1}
Soit
$$
K=\mathbb{Q}(\sqrt{-15})
$$
et, pour tout $a\in K^*$, on considère le polynôme
$$
P(X)=X^3-a\in K[X].
$$
\begin{enumerate}
\item Rappeler (ou démontrer) un critère d'irréductibilité d'un polynôme de degré
      $3$ à coefficients dans un corps $K$.
\item Montrer qu'il existe un élément $a\in K^*$ tel que $P(X)=X^3-a$ soit
      irréductible sur $K$. Dans toute la suite, on fixe un tel élément $a$.
\item Soit $\theta$ une racine de $P$ dans une clôture algébrique de $K$, et soit
      $j=e^{\frac{2i\pi}{3}}$. Montrer que le corps de décomposition de $P$ sur $K$
      est $L=K(\theta,j)=K(\theta,\theta j,\theta j^2)$.
\item Montrer que l'extension $L/K$ est galoisienne de degré $6$ et que
      $\sqrt5\in L$.
\item Montrer qu'il existe des automorphismes $\sigma,\tau\in\operatorname{Gal}(L/K)$
      tels que
      $$
      \sigma(\sqrt5)=\sqrt5,\qquad \sigma(\theta)=j\theta,
      $$
      et
      $$
      \tau(\sqrt5)=-\sqrt5,\qquad \tau(\theta)=\theta.
      $$
\item Déterminer les ordres de $\sigma$ et de $\tau$ dans $\operatorname{Gal}(L/K)$,
      puis calculer $\tau\sigma\tau^{-1}$.
\item En déduire la structure du groupe $\operatorname{Gal}(L/K)$, puis déterminer
      toutes les sous-extensions de $L/K$.
\end{enumerate}

\source{\emph{\small École Normale Supérieure de Kouba (ENS), 2023, Épreuve~n\textsuperscript{o}~03.}}
\end{problem}

\newpage

\begin{problem}{(ENS Kouba) — Épreuve N°03}{alg:ens-kouba-2023-e3-p2}
Soit $K$ un corps algébriquement clos, et soit $A_n=K[X_1,\dots,X_n]$. Pour tout
idéal $I\subset A_n$, on définit
$$
V(I)=\{x\in K^n:\ f(x)=0,\ \forall\,f\in I\},
\qquad
\sqrt I=\{f\in A_n:\ \exists\,m\ge1,\ f^m\in I\}.
$$
De plus, pour toute partie $S\subset K^n$, on pose
$$
I(S)=\{f\in A_n:\ f(x)=0,\ \forall\,x\in S\}.
$$
\begin{enumerate}
\item Montrer que $I(S)$ est un idéal de $A_n$.
\item Montrer que $\sqrt I\subset I(V(I))$.
\item (Astuce de Rabinowitsch) Soit $Q\in I(V(I))$. On plonge $A_n$ dans
      $A_{n+1}=K[X_1,\dots,X_n,X_{n+1}]$, et on considère l'idéal
      $$
      J=\left\langle I\cup\{1-X_{n+1}Q\}\right\rangle\subset A_{n+1}.
      $$
      Montrer que $V(J)=\varnothing$ dans $K^{n+1}$.
\item En déduire, à l'aide du Nullstellensatz faible, que la classe de
      $1-X_{n+1}Q$ est inversible dans $A_{n+1}/\langle I\rangle$.
\item Conclure que $Q\in\sqrt I$, puis en déduire que $I(V(I))=\sqrt I$.
      (On retrouve ainsi le Nullstellensatz de Hilbert.)
\end{enumerate}

\source{\emph{\small École Normale Supérieure de Kouba (ENS), 2023, Épreuve~n\textsuperscript{o}~03.}}
\end{problem}

\newpage

\begin{problem}{(ENS Kouba) — Épreuve N°03}{alg:ens-kouba-2023-e3-p3}
Soit $A$ un anneau commutatif intègre, et soit $a\in A$, avec $a\neq0$.
\begin{enumerate}
\item Étudier l'idéal $(a,X)$ de $A[X]$ et montrer qu'il ne peut être principal que
      si $a$ est inversible dans $A$, c'est-à-dire $a\in A^*$. Dans ce cas, montrer
      que $(a,X)=A[X]$.
\item En déduire que si l'anneau $A[X]$ est principal, alors $A$ est un corps.
\end{enumerate}

\source{\emph{\small École Normale Supérieure de Kouba (ENS), 2023, Épreuve~n\textsuperscript{o}~03.}}
\end{problem}

\newpage

%% ===================================================================
\section{Université Ahmed Ben Bella --- Oran 1}
%% ===================================================================

\begin{problem}{(Oran 1) — Épreuve N°03}{alg:oran1-2023-e3-p2}
Soit $\mathbb{F}_p$ le corps premier à $p$ éléments, $\alpha$ une racine septième de
l'unité et $K = \mathbb{F}_p(\alpha)$.
\begin{enumerate}
\item Expliquer pourquoi $K$ est un corps fini.
\item Supposons que $p = 3$. Déterminer le cardinal du corps $K$.
\item Supposons que $p \neq 3$. Soit $f(x) = x^p - x + 3$ ; si $\beta$ est une racine
      de $f$ dans une extension de $\mathbb{F}_p$, montrer que $f$ admet $p$ racines
      distinctes dans le corps d'extension $\mathbb{F}_p(\beta)$.
\end{enumerate}

\source{\emph{\small Université Ahmed Ben Bella - Oran 1, 2023, Épreuve~n\textsuperscript{o}~03.}}
\end{problem}

\newpage

%% ===================================================================
\section{Université Akli Mohand Oulhadj --- Bouira}
%% ===================================================================

\begin{problem}{(Bouira) — Épreuve N°01}{alg:bouira-2025-e1-p1}
Soit $n > 0$ un nombre entier. On note $\sigma(n)$ la somme des diviseurs de $n$.
Par exemple $\sigma(6) = 1+2+3+6$. On dit que $n$ est parfait si $\sigma(n) = 2n$.
\begin{enumerate}
\item Les nombres $12$ et $28$ sont-ils parfaits ?
\item Soit $n \geq 2$ ; montrer que $\sigma(n) \geq n + 1$.
\item Soient $p_1$ et $p_2$ deux nombres premiers distincts et $\alpha, \beta$ deux
      entiers $\geq 2$.
      \begin{enumerate}
      \item Calculer $\sigma(p_1)$ et $\sigma(p_1^\alpha)$.
      \item Montrer que $\sigma(p_1^\alpha p_2^\beta) = \sigma(p_1^\alpha)\sigma(p_2^\beta)$.
      \item Pour $n = 2025$, calculer $\sigma(n)$.
      \end{enumerate}
\item On suppose que pour tous entiers $n$ et $m$ premiers entre eux,
      $\sigma(nm) = \sigma(n)\sigma(m)$. On pose
      $$
      M_p = 2^{p-1}(2^p - 1),
      $$
      où $p$ et $2^p - 1$ sont premiers.
      \begin{enumerate}
      \item Calculer $\sigma(2^{p-1})$ et $\sigma(2^p - 1)$.
      \item En déduire que $M_p$ est un nombre parfait.
      \end{enumerate}
\item On suppose que $n$ est un nombre parfait pair, de la forme $n = 2^l a$ avec
      $l > 0$ et $a$ impair. Montrer que $2^{l+1} a = (2^{l+1} - 1)\sigma(a)$.
\end{enumerate}

\source{\emph{\small Université Akli Mohand Oulhadj - Bouira, 2025, Épreuve~n\textsuperscript{o}~01.}}
\end{problem}

\newpage

\begin{problem}{(Bouira) — Épreuve N°01}{alg:bouira-2025-e1-p2}
On considère le polynôme
$$
f(x) = x^3 - 7x + 7 \in \mathbb{Q}[x].
$$
\begin{enumerate}
\item Montrer que $f$ est irréductible sur $\mathbb{Q}$.
\item Montrer que $f$ possède trois racines réelles $\alpha, \beta, \gamma$ telles que
      $\alpha > \beta > 0 > \gamma$.
\item En calculant $(\beta - \gamma)^2$ en fonction de $\alpha$, montrer que
      $\beta - \gamma = \dfrac{\alpha}{2\alpha - 3}$.
\end{enumerate}

\source{\emph{\small Université Akli Mohand Oulhadj - Bouira, 2025, Épreuve~n\textsuperscript{o}~01.}}
\end{problem}

\newpage

%% ===================================================================
\section{Université des Sciences et de la Technologie Houari Boumediène (USTHB)}
%% ===================================================================

\begin{problem}{(USTHB) — Épreuve N°03}{alg:usthb-2014-e3-p1}
Les affirmations suivantes sont-elles vraies ? Justifier.
\begin{enumerate}
\item Tout code linéaire est cyclique.
\item Un anneau à chaîne est un anneau de Galois.
\item Tout anneau principal est à chaîne.
\item La caractéristique d'un anneau de Galois est toujours première.
\item Un anneau de Galois n'a pas d'éléments nilpotents.
\item Tout code libre est auto-dual.
\end{enumerate}

\source{\emph{\small Université des Sciences et de la Technologie Houari Boumediène (USTHB), 2014, Épreuve~n\textsuperscript{o}~03.}}
\end{problem}

\newpage

\begin{problem}{(USTHB) — Épreuve N°03}{alg:usthb-2014-e3-p2}
Soit $(n,e)$ une clé publique RSA.
\begin{enumerate}
\item Soient $p,q$ deux premiers distincts avec $p\equiv q\equiv 2 \pmod 3$,
      $n=pq$ et
      $$
      e=\frac{2(p-1)(q-1)+1}{3}.
      $$
      Montrer que $e$ est premier avec $\varphi(n)$ et calculer son inverse modulo
      $\varphi(n)$.
\item On prend $(n,e)=(187,107)$ et le cryptogramme reçu est $9$. Quel est le
      message ?
\end{enumerate}

\source{\emph{\small Université des Sciences et de la Technologie Houari Boumediène (USTHB), 2014, Épreuve~n\textsuperscript{o}~03.}}
\end{problem}

\newpage

\begin{problem}{(USTHB) — Épreuve N°03}{alg:usthb-2014-e3-p3}
Les matrices suivantes peuvent-elles être génératrices de codes sur $\mathbb{Z}_4$ ?
$$
G_1=\begin{pmatrix} 1&1&1&0\\ 0&0&2&2\\ 0&2&0&2\end{pmatrix},
\qquad
G_2=\begin{pmatrix} 1&0&0&0&5&1&2&1\\ 0&1&0&0&1&2&7&1\\ 0&0&1&0&3&6&3&2\\ 0&0&0&1&2&3&1&1\end{pmatrix}.
$$
Si oui :
\begin{enumerate}
\item déterminer le type et la longueur ;
\item le code est-il auto-dual ? Donner la matrice de contrôle ;
\item le code est-il libre ?
\end{enumerate}

\source{\emph{\small Université des Sciences et de la Technologie Houari Boumediène (USTHB), 2014, Épreuve~n\textsuperscript{o}~03.}}
\end{problem}

\newpage

\begin{problem}{(USTHB) — Épreuve N°03}{alg:usthb-2014-e3-p4}
Déterminer tous les anneaux de Galois d'ordre $64$, avec leurs caractéristiques,
leurs corps résiduels et leurs groupes multiplicatifs.

\source{\emph{\small Université des Sciences et de la Technologie Houari Boumediène (USTHB), 2014, Épreuve~n\textsuperscript{o}~03.}}
\end{problem}

\newpage

\begin{problem}{(USTHB) — Épreuve N°04}{alg:usthb-2014-e4-p1}
\begin{enumerate}
\item Soient $H$ et $K$ deux sous-groupes distingués et résolubles d'un groupe $G$.
      Montrer que $HK$ est distingué et résoluble.
\item Soient $H,K$ distingués tels que $G/H$ et $G/K$ soient résolubles. Montrer
      que $G/(H\cap K)$ est résoluble.
\item Déterminer à isomorphisme près les groupes d'ordre $15$ et d'ordre $35$.
\item Déterminer une suite de Jordan--Hölder pour :
      \begin{enumerate}
      \item $\mathbb{Z}/18\mathbb{Z}$,
      \item $\mathbb{Z}/2\mathbb{Z}\times\mathbb{Z}/2\mathbb{Z}$,
      \item $S_3$, $S_4$, $S_5$.
      \end{enumerate}
\end{enumerate}

\source{\emph{\small Université des Sciences et de la Technologie Houari Boumediène (USTHB), 2014, Épreuve~n\textsuperscript{o}~04.}}
\end{problem}

\newpage

\begin{problem}{(USTHB) — Épreuve N°04}{alg:usthb-2014-e4-p2}
Soit $A$ un anneau.
\begin{enumerate}
\item Montrer par un contre-exemple que l'ensemble des éléments nilpotents de $A$
      ne forme pas un sous-groupe abélien de $A$ (on pourra choisir
      $A=M_2(\mathbb{C})$).
\item Soit $N=\{a\in A \ / \ ax \text{ est nilpotent } \forall x\in A\}$. Montrer que
      $N$ est un idéal bilatère de $A$ et que tout élément de $N$ est nilpotent.
\item Soit $I$ un idéal bilatère de $A$ dont tout élément est nilpotent. Montrer que
      $I\subset N$.
\end{enumerate}

\source{\emph{\small Université des Sciences et de la Technologie Houari Boumediène (USTHB), 2014, Épreuve~n\textsuperscript{o}~04.}}
\end{problem}

\newpage

\begin{problem}{(USTHB) — Épreuve N°04}{alg:usthb-2014-e4-p3}
Soit $K$ un corps. On pose
$$
A=\frac{K[X,Y]}{(X^2,XY,Y^2)}.
$$
\begin{enumerate}
\item Déterminer les éléments inversibles de $A$.
\item Déterminer tous les idéaux principaux de $A$.
\item Déterminer tous les idéaux de $A$.
\end{enumerate}

\source{\emph{\small Université des Sciences et de la Technologie Houari Boumediène (USTHB), 2014, Épreuve~n\textsuperscript{o}~04.}}
\end{problem}

\newpage

\begin{problem}{(USTHB) — Épreuve N°04}{alg:usthb-2014-e4-p4}
Soient $L,M$ deux $A$-modules et $f:L\to M$ un homomorphisme.
\begin{enumerate}
\item On suppose que $\operatorname{Ker}(f)$ et $\operatorname{Im}(f)$ sont de type
      fini. Montrer que $L$ est de type fini.
\item On suppose que $\operatorname{Ker}(f)\cong A^p$ et
      $\operatorname{Im}(f)\cong A^q$. Montrer que $L\cong A^{p+q}$.
\end{enumerate}

\source{\emph{\small Université des Sciences et de la Technologie Houari Boumediène (USTHB), 2014, Épreuve~n\textsuperscript{o}~04.}}
\end{problem}

\newpage

\begin{problem}{(USTHB) — Épreuve N°02}{alg:usthb-2015-e2-p1}
Soit $A = \mathbb{C}^{\mathbb{N}^*}$ l'ensemble des fonctions arithmétiques muni de
l'addition usuelle et du produit de Dirichlet
$$
(f*g)(n)=\sum_{d\mid n} f(d)\,g(n/d).
$$
\begin{enumerate}
\item Préciser l'élément neutre $\delta$ de la multiplication.
\item Montrer que le groupe des unités de $A$ est $U(A)=\{f\in A : f(1)\neq 0\}$.
\item On dit qu'une fonction arithmétique $f$ est multiplicative si $f(1)\neq 0$ et
      $(n,m)=1 \Rightarrow f(nm)=f(n)f(m)$. On désigne par $\mathcal{M}$ l'ensemble
      des fonctions arithmétiques multiplicatives. Montrer que $\mathcal{M}$ est un
      sous-groupe de $U(A)$.
\item On désigne par $1$ et $j$ les fonctions définies par $1(n)=1$ et $j(n)=n$.
      \begin{enumerate}
      \item Montrer que $1$ et $j$ appartiennent à $\mathcal{M}$.
      \item On désigne par $\mu$ la fonction de Möbius et par $\varphi$ la fonction
            indicatrice d'Euler. Prouver que $\mu$ est l'inverse pour le produit de
            Dirichlet de la fonction arithmétique $1$ et que $1*\varphi=j$.
      \end{enumerate}
\end{enumerate}

\source{\emph{\small Université des Sciences et de la Technologie Houari Boumediène (USTHB), 2015, Épreuve~n\textsuperscript{o}~02.}}
\end{problem}

\newpage

\begin{problem}{(USTHB) — Épreuve N°01}{alg:usthb-2023-e1-p1}
Soit $D$ l'ensemble des nombres décimaux, c'est-à-dire
$$
D=\mathbb{Z}\Big[\tfrac1{10}\Big]=\Big\{\tfrac{a}{10^n}\ :\ a\in\mathbb{Z},\ n\in\mathbb{N}\Big\},
$$
qui est un sous-anneau de $\mathbb{Q}$. Soit $I$ un idéal de $D$.
\begin{enumerate}
\item Montrer que $I\cap\mathbb{Z}$ est un idéal de $\mathbb{Z}$.
\item En déduire que $I$ est un idéal principal de $D$.
\end{enumerate}

\source{\emph{\small Université des Sciences et de la Technologie Houari Boumediène (USTHB), 2023, Épreuve~n\textsuperscript{o}~01.}}
\end{problem}

\newpage

\begin{problem}{(USTHB) — Épreuve N°01}{alg:usthb-2023-e1-p2}
Soit $A$ un anneau et soit
$$
0\longrightarrow M\xrightarrow{\ f\ }N\xrightarrow{\ g\ }P\longrightarrow0
$$
une suite exacte de $A$-modules. On suppose que $M$ et $P$ sont de type fini.
Montrer que $N$ est de type fini.

\source{\emph{\small Université des Sciences et de la Technologie Houari Boumediène (USTHB), 2023, Épreuve~n\textsuperscript{o}~01.}}
\end{problem}

\newpage

\begin{problem}{(USTHB) — Épreuve N°01}{alg:usthb-2023-e1-p3}
Soit $n\ge2$ un entier. Étudier l'exactitude des deux suites de
$\mathbb{Z}$-modules suivantes.
\begin{enumerate}
\item[\textbf{(S1)}] $0\longrightarrow\mathbb{Z}\xrightarrow{\ \times n\ }\mathbb{Z}\xrightarrow{\ \pi\ }\mathbb{Z}/n\mathbb{Z}\longrightarrow0$,
      où $\pi$ est la projection canonique.
\item[\textbf{(S2)}] $0\longrightarrow\mathbb{Z}/n\mathbb{Z}\xrightarrow{\ f\ }\mathbb{Z}/n^2\mathbb{Z}\xrightarrow{\ g\ }\mathbb{Z}/n\mathbb{Z}\longrightarrow0$,
      où $f(\overline x)=\overline{nx}$ et $g$ est la réduction modulo $n$.
\end{enumerate}

\source{\emph{\small Université des Sciences et de la Technologie Houari Boumediène (USTHB), 2023, Épreuve~n\textsuperscript{o}~01.}}
\end{problem}

\newpage

\begin{problem}{(USTHB) — Épreuve N°01}{alg:usthb-2023-e1-p4}
Soit $f=x^5-x+1\in\mathbb{Q}[x]$, $\alpha$ une racine de $f$ et
$K=\mathbb{Q}(\alpha)$. On rappelle que pour le trinôme $x^n+ax+b$,
$$
\operatorname{disc}=(-1)^{n(n-1)/2}\big[(-1)^{n-1}(n-1)^{n-1}a^n+n^nb^{n-1}\big].
$$
\begin{enumerate}
\item Montrer que $f$ est irréductible sur $\mathbb{Q}$.
\item Calculer le discriminant de $f$, puis déterminer l'anneau des entiers
      $\mathcal{O}_K$ et le discriminant $\operatorname{disc}(K)$.
\item Décomposer les idéaux $3\,\mathcal{O}_K$ et $19\,\mathcal{O}_K$ en produit
      d'idéaux premiers.
\item Déterminer le nombre de racines réelles de $f$, puis montrer que le groupe de
      Galois de $f$ sur $\mathbb{Q}$ est $G\cong S_5$.
\end{enumerate}

\source{\emph{\small Université des Sciences et de la Technologie Houari Boumediène (USTHB), 2023, Épreuve~n\textsuperscript{o}~01.}}
\end{problem}

\newpage

\begin{problem}{(USTHB) — Épreuve N°02}{alg:usthb-2026-e2-p1}
On note $A$ l'ensemble des éléments de $\mathbb{C}$ qui sont entiers sur
$\mathbb{Z}$. Soient $\alpha\in A\setminus\{0\}$ et $f(X)$ son polynôme minimal sur
$\mathbb{Q}$.

Les équivalences suivantes sont-elles vraies ?
$$
\frac1\alpha\in A\Longleftrightarrow f(0)=\pm1\Longleftrightarrow\frac1\alpha\in\mathbb{Z}[\alpha].
$$
Justifier votre réponse.

\source{\emph{\small Université des Sciences et de la Technologie Houari Boumediène (USTHB), 2026, Épreuve~n\textsuperscript{o}~02 — barème : 4 points.}}
\end{problem}

\newpage

\begin{problem}{(USTHB) — Épreuve N°02}{alg:usthb-2026-e2-p2}
Soit le polynôme $f(X)=X^3+X^2-2X+8$ et $\alpha\in\mathbb{C}$ tel que
$f(\alpha)=0$. On pose $\beta=4/\alpha$ et $K=\mathbb{Q}(\beta)$. On note $A$
l'anneau des entiers de $K$ sur $\mathbb{Z}$.
\begin{enumerate}
\item Donner $D_{K/\mathbb{Q}}(1,\alpha,\alpha^2)$ et
      $D_{K/\mathbb{Q}}(1,\alpha,\beta)$.
\item Est-ce que $(1,\alpha,\alpha^2)$ est une base de $A$ sur $\mathbb{Z}$ ?
\end{enumerate}
Indication : on pourra utiliser le fait que $503$ est un nombre premier.

\source{\emph{\small Université des Sciences et de la Technologie Houari Boumediène (USTHB), 2026, Épreuve~n\textsuperscript{o}~02 — barème : 5 points.}}
\end{problem}

\newpage

\begin{problem}{(USTHB) — Épreuve N°02}{alg:usthb-2026-e2-p3}
\begin{enumerate}
\item Soient $M$ et $N$ deux $A$-modules noethériens. Montrer que $M\times N$ est
      noethérien.
\item Soient $E$ un $A$-module, $M$ et $N$ deux sous-modules de $E$. Montrer que si
      $M$ et $N$ sont noethériens alors $M+N$ est noethérien.
\item Soit $M$ un module sur un anneau noethérien. Si $M$ est de type fini alors
      $M$ est noethérien.
\item Soit $A$ un anneau et $S$ une partie multiplicative de $A$. Si $A$ est
      noethérien alors $S^{-1}A$ est noethérien.
\end{enumerate}

\source{\emph{\small Université des Sciences et de la Technologie Houari Boumediène (USTHB), 2026, Épreuve~n\textsuperscript{o}~02 — barème : 5 points.}}
\end{problem}

\newpage

\begin{problem}{(USTHB) — Épreuve N°02}{alg:usthb-2026-e2-p4}
Soit $n$ un entier $\ge3$.
\begin{enumerate}
\item Montrer que le degré de l'extension
      $\mathbb{Q}(\cos(2\pi/n))/\mathbb{Q}$ est $\varphi(n)/2$, où $\varphi$ désigne
      la fonction indicatrice d'Euler.
\item Soit $x\in\mathbb{Q}$ écrit sous forme irréductible $a/b$, avec $a$ et $b$
      entiers premiers entre eux et $b>0$. On suppose $b\ge3$. Calculer le degré de
      $\mathbb{Q}(\cos(2\pi a/b))/\mathbb{Q}$.
\item Pour $n\ge5$, en déduire que le degré de
      $\mathbb{Q}(\sin(2\pi/n))/\mathbb{Q}$ est
      $$
      \begin{cases}
      \varphi(n), & \text{si } \operatorname{pgcd}(n,8)<4,\\
      \dfrac{\varphi(n)}4, & \text{si } \operatorname{pgcd}(n,8)=4,\\
      \dfrac{\varphi(n)}2, & \text{si } \operatorname{pgcd}(n,8)>4.
      \end{cases}
      $$
\item Application : déterminer le polynôme minimal de $\sin(2\pi/5)$ sur
      $\mathbb{Q}$.
\end{enumerate}
Indication : $\sin(5x)=16\sin(x)^5-20\sin(x)^3+5\sin(x)$ pour tout $x$ réel.

\source{\emph{\small Université des Sciences et de la Technologie Houari Boumediène (USTHB), 2026, Épreuve~n\textsuperscript{o}~02 — barème : 6 points.}}
\end{problem}

\newpage

%% ===================================================================
\section{Université Ferhat Abbas --- Sétif 1}
%% ===================================================================

\begin{problem}{(Sétif 1) — Épreuve N°01}{alg:setif1-2025-e1-p1}
Soit $G$ un groupe fini simple non abélien, et $M$ un sous-groupe maximal propre de
$G$ (i.e. $M\ne G$ et il n'existe aucun sous-groupe strictement compris entre $M$
et $G$).
\begin{enumerate}
\item Montrer que le normalisateur $N_G(M)$ de $M$ dans $G$ est égal à $M$.
\item En déduire que l'action de $G$ par translation sur l'ensemble des classes à
      gauche $G/M$ est fidèle et transitive, et que $|G|$ divise $[G:M]!$.
\item Application : montrer que si $G$ simple non abélien possède un sous-groupe
      maximal $M$ d'indice $n$, alors $G$ se plonge dans $S_n$ (groupe symétrique),
      et donc $|G|$ divise $n!$.
\end{enumerate}

\source{\emph{\small Université Ferhat Abbas - Sétif 1, 2025, Épreuve~n\textsuperscript{o}~01.}}
\end{problem}

\newpage

\begin{problem}{(Sétif 1) — Épreuve N°01}{alg:setif1-2025-e1-p2}
Soit $G$ un groupe fini. On rappelle que le sous-groupe de Frattini
$\mathrm{Frat}(G)$ (ou $\Phi(G)$) est l'intersection de tous les sous-groupes
maximaux de $G$.
\begin{enumerate}
\item Montrer que $\mathrm{Frat}(G)$ est un sous-groupe caractéristique (donc
      normal) de $G$.
\item Montrer que si $x\in G$ est tel que $G=\langle \mathrm{Frat}(G),x\rangle$,
      alors $G=\langle x\rangle$ (propriété dite de « non-générateur »).
\item Montrer que pour $G$ un $p$-groupe fini,
      $\mathrm{Frat}(G)=G'G^p$, sous-groupe engendré par le dérivé et les puissances
      $p$-ièmes.
\end{enumerate}

\source{\emph{\small Université Ferhat Abbas - Sétif 1, 2025, Épreuve~n\textsuperscript{o}~01.}}
\end{problem}

\newpage

\begin{problem}{(Sétif 1) — Épreuve N°01}{alg:setif1-2025-e1-p3}
Soit $G$ un groupe d'ordre $p^3q$ où $p,q$ sont des nombres premiers distincts. On
veut montrer que $G$ n'est pas simple (sauf cas triviaux à exclure).
\begin{enumerate}
\item Rappeler les théorèmes de Sylow (existence, conjugaison, nombre de $p$-Sylow
      $n_p\equiv1\pmod p$ et $n_p\mid q$ ; nombre de $q$-Sylow
      $n_q\equiv1\pmod q$ et $n_q\mid p^3$).
\item Traiter le cas $q<p$ : montrer que $n_p=1$, donc que le $p$-Sylow est normal
      et que $G$ n'est pas simple.
\item Traiter le cas $q>p$ : discuter les valeurs possibles de $n_q$ et montrer que
      dans les configurations restantes, $G$ n'est pas simple (on pourra utiliser un
      argument de comptage d'éléments).
\end{enumerate}

\source{\emph{\small Université Ferhat Abbas - Sétif 1, 2025, Épreuve~n\textsuperscript{o}~01.}}
\end{problem}

\newpage

\begin{problem}{(Sétif 1) — Épreuve N°02}{alg:setif1-2025-e2-p1}
Soit $G$ un groupe et $M$ un sous-groupe maximal de $G$, simple et distingué dans
$G$. Soit $N$ un sous-groupe distingué de $G$ non trivial ($N\neq\{e\}$ et
$N\neq G$).
\begin{enumerate}
\item Montrer que si $N\subseteq M$, alors $N=M$.
\item Montrer que si $N\not\subseteq M$, alors $NM=G$ et $G/N$ est isomorphe à un
      quotient de $M$ ; en déduire que $G/N$ est simple ou trivial.
\item Conclure : tout quotient propre non trivial de $G$ est simple.
\end{enumerate}

\source{\emph{\small Université Ferhat Abbas - Sétif 1, 2025, Épreuve~n\textsuperscript{o}~02 — barème : 6 points. Le même sujet circule aussi sans corrigé sous un autre numéro d'épreuve.}}
\end{problem}

\newpage

\begin{problem}{(Sétif 1) — Épreuve N°02}{alg:setif1-2025-e2-p2}
Soit $G$ un groupe fini et $\Phi(G)$ son sous-groupe de Frattini (intersection des
sous-groupes maximaux de $G$). Montrer l'équivalence des assertions suivantes :
\begin{enumerate}
\item $G$ est nilpotent ;
\item $G/\Phi(G)$ est nilpotent ;
\item $G'\subseteq\Phi(G)$ (où $G'$ est le groupe dérivé), i.e. $G/\Phi(G)$ est
      abélien.
\end{enumerate}

\source{\emph{\small Université Ferhat Abbas - Sétif 1, 2025, Épreuve~n\textsuperscript{o}~02 — barème : 6 points.}}
\end{problem}

\newpage

\begin{problem}{(Sétif 1) — Épreuve N°02}{alg:setif1-2025-e2-p3}
Soient $p$ et $q$ deux nombres premiers distincts et $G$ un groupe d'ordre $p^3q$.
On veut montrer que $G$ possède un sous-groupe de Sylow distingué, sauf exception à
préciser.
\begin{enumerate}
\item Rappeler les théorèmes de Sylow sur $n_p$ et $n_q$ (nombres de $p$- et
      $q$-Sylow).
\item Montrer que si $n_p=q$, alors $q\equiv 1 \pmod p$ et donc $q>p$.
\item On suppose $n_p\neq 1$ et $n_q\neq 1$. Montrer que $n_q\in\{p,p^2,p^3\}$,
      éliminer le cas $n_q=p$, puis montrer que le cas $n_q=p^3$ conduit à une
      contradiction.
\item Étudier le cas restant $n_q=p^2$ et conclure (on examinera la petite exception
      numérique).
\end{enumerate}

\source{\emph{\small Université Ferhat Abbas - Sétif 1, 2025, Épreuve~n\textsuperscript{o}~02 — barème : 8 points.}}
\end{problem}

\newpage

%% ===================================================================
\section{Université Kasdi Merbah --- Ouargla}
%% ===================================================================

\begin{problem}{(Ouargla) — Épreuve N°35}{alg:ouargla-2015-e35-p1}
Soit $A$ un sous-anneau de $\mathbb{Z}$.
\begin{enumerate}
\item Montrer que $A$ est un sous-anneau principal de la forme $n\mathbb{Z}$ pour un
      certain entier $n$.
\item Soit $x\in A$. On pose $a=x+1$ avec $a\in A$ ; vérifier les propriétés
      indiquées dans l'énoncé original.
\item Montrer que $\mathbb{Z}$ est commutatif.
\end{enumerate}

\source{\emph{\small Université Kasdi Merbah - Ouargla, 2015, Épreuve~n\textsuperscript{o}~35 — barème : 8 points. Transcription partielle : la photo d'origine est peu lisible, la question~2 doit être recoupée avec le sujet papier.}}
\end{problem}

\newpage

\begin{problem}{(Ouargla) — Épreuve commune}{alg:ouargla-2022-gen-p1}
Soit $A$ un anneau d'élément neutre $0$ et d'élément unité $1$. On note $Z(A)$ le
centre de $A$ et l'on suppose que
$$
\forall x\in A,\qquad x^2-x\in Z(A).
$$
\begin{enumerate}
\item Montrer que $Z(A)$ est un sous-anneau de $A$.
\item Soient $x,y\in A$. On pose $\alpha=x+y$.
      \begin{enumerate}
      \item Vérifier que $\alpha^2-\alpha\in Z(A)$, puis en déduire que
            $xy+yx\in Z(A)$.
      \item En utilisant la question précédente, montrer que $x^2\in Z(A)$.
      \item Montrer que $x\in Z(A)$.
      \end{enumerate}
\item En déduire que l'anneau $A$ est commutatif.
\end{enumerate}

\source{\emph{\small Université Kasdi Merbah - Ouargla, 2022, épreuve commune — barème : 8 points.}}
\end{problem}

\newpage

\begin{problem}{(Ouargla) — Épreuve commune}{alg:ouargla-2022-gen-p2}
\begin{enumerate}
\item Soit $n\in\mathbb{N}^*$ et
      $$
      P(x)=a_0+a_1x+\cdots+a_nx^n\in\mathbb{Z}[X],\qquad a_n\neq0.
      $$
      Soit $\alpha=\dfrac{a}{b}\in\mathbb{Q}$, avec $\gcd(a,b)=1$ et $b\neq0$.
      Montrer que si $\alpha$ est une racine de $P$, alors $a\mid a_0$ et
      $b\mid a_n$. En déduire que si $P$ est unitaire, alors toute racine
      rationnelle est un entier.
\item On considère le polynôme $P(x)=x^4-3x^3+11x^2-9x$. Déterminer :
      \begin{enumerate}
      \item toutes les racines de $P$ dans $\mathbb{C}$ ;
      \item le corps de décomposition $E$ de $P$ sur $\mathbb{Q}$.
      \end{enumerate}
\end{enumerate}

\source{\emph{\small Université Kasdi Merbah - Ouargla, 2022, épreuve commune — barème : 8 points.}}
\end{problem}

\newpage

\begin{problem}{(Ouargla) — Épreuve commune}{alg:ouargla-2022-gen-p3}
On rappelle que si $G$ est un groupe fini et $a,b\in G$ tels que $o(a)=n$,
$o(b)=m$ et $\langle a\rangle\cap\langle b\rangle=\{e\}$, alors
$o(ab)=\operatorname{ppcm}(n,m)$.

Soit $G$ un groupe fini d'ordre $35$.
\begin{enumerate}
\item Déterminer $n_5$ et $n_7$, le nombre de $5$-sous-groupes de Sylow et de
      $7$-sous-groupes de Sylow de $G$, respectivement.
\item En déduire que $G$ est cyclique.
\end{enumerate}
Soit $E=\mathbb{Q}\!\left(\sqrt{2},\,i\right)$.
\begin{enumerate}
\setcounter{enumi}{2}
\item Déterminer le degré $[E:\mathbb{Q}]$.
\item Déterminer le groupe de Galois $G=\operatorname{Gal}(E/\mathbb{Q})$.
\item Déterminer tous les sous-groupes de $G$ ainsi que tous les sous-corps
      intermédiaires de $E$.
\item Déterminer un élément primitif $\theta$ de $E$ ainsi que son inverse.
\item Donner deux bases différentes de $E$ sur $\mathbb{Q}$.
\end{enumerate}

\source{\emph{\small Université Kasdi Merbah - Ouargla, 2022, épreuve commune — barème : 4 points.}}
\end{problem}

\newpage

\begin{problem}{(Ouargla) — Épreuve de spécialité}{alg:ouargla-2022-spec-p3}
On rappelle que si $G$ est un groupe fini et $a,b\in G$ tels que $o(a)=n$ et
$o(b)=m$, avec $\langle a\rangle\cap\langle b\rangle=\{e\}$, alors
$o(ab)=\operatorname{ppcm}(n,m)$.

Soit $G$ un groupe fini d'ordre $35$.
\begin{enumerate}
\item Déterminer $n_5$ et $n_7$, le nombre de $5$-sous-groupes de Sylow et le nombre
      de $7$-sous-groupes de Sylow de $G$, respectivement.
\item En déduire que $G$ est cyclique.
\end{enumerate}
Soit $E=\mathbb{Q}\!\left(\sqrt{2},\sqrt{3}\right)$.
\begin{enumerate}
\setcounter{enumi}{2}
\item Déterminer le degré $[E:\mathbb{Q}]$.
\item Déterminer le groupe de Galois $G=\operatorname{Gal}(E/\mathbb{Q})$.
\item Déterminer tous les sous-groupes de $G$ ainsi que les sous-corps
      intermédiaires de $E$.
\item Déterminer un élément primitif $\theta$ de $E$ ainsi que son inverse.
\item Donner deux bases différentes de $E$ sur $\mathbb{Q}$.
\end{enumerate}

\source{\emph{\small Université Kasdi Merbah - Ouargla, 2022, épreuve de spécialité — barème : 4 points. Variante de l'exercice précédent : seul le corps $E$ change ($\sqrt3$ au lieu de $i$).}}
\end{problem}

\newpage

%% ===================================================================
\section{Université Mohamed Boudiaf --- M'Sila}
%% ===================================================================

\begin{problem}{(M'Sila) — Épreuve N°34}{alg:msila-2015-e34-p1}
Soit $R$ une relation d'équivalence définie sur un ensemble $E$. Une partie
$X\subset\mathcal P(E)$ est dite saturée par $R$ si
$\bigcup_{x\in X}\overline{x}=X$, avec $\overline{x}$ la classe de $x$ modulo $R$.

Soit $f:E\to F$ une application et soit $R_f$ l'équivalence associée dans $E$,
définie par $x\,R_f\,y\Longleftrightarrow f(x)=f(y)$.
\begin{enumerate}
\item À quelle condition une partie $X\subset E$ est-elle saturée par $R_f$ ?
\item Soit $B$ une partie quelconque de $E$ ; déterminer la plus petite partie
      saturée $\overline B$ contenant $B$.
\end{enumerate}

\source{\emph{\small Université Mohamed Boudiaf - M'Sila, 2015, Épreuve~n\textsuperscript{o}~34 — barème : 5 points.}}
\end{problem}

\newpage

\begin{problem}{(M'Sila) — Épreuve N°11}{alg:msila-2016-e11-p3}
Soit $\mathbb{F}_p^*$ le groupe multiplicatif des éléments non nuls du corps fini
$\mathbb{F}_p=\mathbb{Z}/p\mathbb{Z}$, $p$ étant un nombre premier différent de $2$.
On désigne par $G_p$ l'ensemble des carrés dans $\mathbb{F}_p^*$, c'est-à-dire
$G_p = \{y \in \mathbb{F}_p^* \ / \ \exists x \in \mathbb{F}_p^* : y = x^2\}$.
\begin{enumerate}
\item Vérifier que $G_p$ est un sous-groupe de $\mathbb{F}_p^*$. Montrer que
      l'application $\varphi : x \mapsto x^2$ est un morphisme du groupe
      $\mathbb{F}_p^*$ sur le groupe $G_p$.
\item Calculer $\operatorname{Ker}\varphi$. En déduire que $G_p$ est isomorphe à
      $\mathbb{F}_p^*/\operatorname{Ker}\varphi$.
\item Montrer que $|G_p| = (p - 1)/2$. Pour $p = 7$, calculer $G_p$.
\end{enumerate}

\source{\emph{\small Université Mohamed Boudiaf - M'Sila, 2016, Épreuve~n\textsuperscript{o}~11.}}
\end{problem}

\newpage

\begin{problem}{(M'Sila) — Épreuve N°12}{alg:msila-2016-e12-p1}
Soit $E$ un ensemble à $n$ éléments. Quel est le nombre des couples $(X, Y)$ de
$\mathcal{P}(E) \times \mathcal{P}(E)$ tels que $X \subset Y$ ?

\source{\emph{\small Université Mohamed Boudiaf - M'Sila, 2016, Épreuve~n\textsuperscript{o}~12 — barème : 5 points.}}
\end{problem}

\newpage

\begin{problem}{(M'Sila) — Épreuve N°12}{alg:msila-2016-e12-p2}
\begin{enumerate}
\item Montrer que le polynôme $p(x) = x^4 + x + 1$ est irréductible sur le corps
      fini $\mathbb{F}_2$.
\item Construire $\mathbb{F}_{16}$.
\item Soit $\alpha$ un élément primitif de $\mathbb{F}_{16}$. Montrer que
      $\mathbb{F}_4 = \{0, 1, \alpha^5, \alpha^{10}\}$.
\end{enumerate}

\source{\emph{\small Université Mohamed Boudiaf - M'Sila, 2016, Épreuve~n\textsuperscript{o}~12 — barème : 7 points. Cet exercice est reposé à l'identique dans plusieurs sujets de M'Sila (2015 et 2016).}}
\end{problem}

\newpage

\begin{problem}{(M'Sila) — Épreuve N°14}{alg:msila-2016-e14-p1}
Soit $(P, \le)$ un ensemble partiellement ordonné et $X \subseteq P$, non vide.
\begin{enumerate}
\item Montrer que $\vee X$ et $\wedge X$, s'ils existent, sont uniques.
\item Montrer que si $X = \emptyset$, alors $\vee X$ et $\wedge X$ existent si et
      seulement si $P$ est fermé.
\end{enumerate}

\source{\emph{\small Université Mohamed Boudiaf - M'Sila, 2016, Épreuve~n\textsuperscript{o}~14 — barème : 5 points.}}
\end{problem}

\newpage

\begin{problem}{(M'Sila) — Épreuve N°14}{alg:msila-2016-e14-p2}
Montrer que dans un ensemble ordonné $(X, \le)$ on a
$$
x \ne \sup\, x\ \text{ si et seulement si } \ x \ne \sup A, \quad \forall A \subset X \ \text{tel que } x \in A.
$$
\emph{Indication} : $x = \{a \in X \ / \ a < x\}$.

\source{\emph{\small Université Mohamed Boudiaf - M'Sila, 2016, Épreuve~n\textsuperscript{o}~14 — barème : 4 points. L'énoncé original est ambigu sur la notation $\sup x$ ; il est reproduit tel quel.}}
\end{problem}

\newpage

\begin{problem}{(M'Sila) — Épreuve N°14}{alg:msila-2016-e14-p3}
Soit $E = \{a, b, c, d, e\}$.
\begin{enumerate}
\item Trouver le nombre des relations réflexives sur $E$.
\item Trouver le nombre des relations irréflexives sur $E$.
\item Trouver le nombre des relations symétriques sur $E$.
\item Montrer qu'il existe une relation d'équivalence sur $E$ qui soit incluse dans
      toute autre relation d'équivalence sur $E$.
\item Montrer qu'il existe une relation d'équivalence sur $E$ qui contient toute
      autre relation d'équivalence sur $E$.
\item Supposons que $E$ est un treillis. Donner tous les diagrammes de Hasse de $E$
      (à isomorphisme près).
\end{enumerate}

\source{\emph{\small Université Mohamed Boudiaf - M'Sila, 2016, Épreuve~n\textsuperscript{o}~14 — barème : 11 points. Cet exercice est reposé à l'identique dans plusieurs sujets de M'Sila (2015 et 2016).}}
\end{problem}

\newpage

\begin{problem}{(M'Sila) — Épreuve N°15}{alg:msila-2016-e15-p2}
Soit $E = \{a, b, c\}$ et $\delta$ la relation définie sur $E$ par la matrice
$$
\begin{pmatrix}1&0&1\\0&1&1\\0&0&1\end{pmatrix},
$$
les lignes et les colonnes étant indexées par $a, b, c$ dans cet ordre.
\begin{enumerate}
\item Montrer que $\delta$ est une relation d'ordre.
\item Donner la dimension de $\delta$.
\end{enumerate}

\source{\emph{\small Université Mohamed Boudiaf - M'Sila, 2016, Épreuve~n\textsuperscript{o}~15 — barème : 3 points.}}
\end{problem}

\newpage

\begin{problem}{(M'Sila) — Épreuve N°15}{alg:msila-2016-e15-p3}
\begin{enumerate}
\item Montrer que dans un treillis $T$ on a
      $$
      \begin{cases} x \wedge y = x \wedge z \\ x \vee y = x \vee z \end{cases}
      \Longrightarrow y = z
      $$
      si et seulement si $T$ est distributif.
\item Montrer que l'image d'un treillis distributif par un morphisme de treillis est
      un treillis distributif.
\end{enumerate}

\source{\emph{\small Université Mohamed Boudiaf - M'Sila, 2016, Épreuve~n\textsuperscript{o}~15 — barème : 6 points.}}
\end{problem}

\newpage

\begin{problem}{(M'Sila) — Épreuve N°16}{alg:msila-2016-e16-p1}
Soit $(X, \le)$ un ensemble ordonné, et $E$ une famille de sous-ensembles de $X$
(i.e. $E \subseteq \mathcal{P}(X)$). La relation $\mathfrak{R}$ est définie sur $E$
par :
$$
A\,\mathfrak{R}\,B \ \text{ssi} \
\begin{cases}
\text{pour tout } a \in A \text{ il existe } b \in B \text{ tel que } a \le b, \\
\text{pour tout } b \in B \text{ il existe } a \in A \text{ tel que } a \le b.
\end{cases}
$$
Est-elle réflexive ? symétrique ? antisymétrique ? transitive ?

\source{\emph{\small Université Mohamed Boudiaf - M'Sila, 2016, Épreuve~n\textsuperscript{o}~16 — barème : 6 points.}}
\end{problem}

\newpage

\begin{problem}{(M'Sila) — Épreuve N°16}{alg:msila-2016-e16-p3}
Soit $E = \{x, y, z, t\}$ et $R$ la relation définie sur $E$ par la matrice
$$
\begin{pmatrix}1&1&1&1\\0&1&0&1\\0&0&1&1\\0&0&0&1\end{pmatrix},
$$
les lignes et les colonnes étant indexées par $x, y, z, t$ dans cet ordre.
\begin{enumerate}
\item Montrer que $R$ est une relation d'ordre.
\item Donner la dimension de $R$.
\end{enumerate}

\source{\emph{\small Université Mohamed Boudiaf - M'Sila, 2016, Épreuve~n\textsuperscript{o}~16 — barème : 4 points.}}
\end{problem}

\newpage

%% ===================================================================
\section{Université Mohamed Seddik Benyahia --- Jijel}
%% ===================================================================

\begin{problem}{(Jijel) — Épreuve N°03}{alg:jijel-2019-e3-p2}
Soit $A = \{a_1, -a_2\}$ et $E = \{p_1, -p_2\}$ deux alphabets. On donne une
identité de fonctions génératrices reliant les fonctions symétriques complètes
$h_n$ des alphabets $A$ et $E$ à des séries rationnelles en $z$.
\begin{enumerate}
\item Déterminer la fonction génératrice de $\displaystyle\sum_{n \ge 0} F_{n+1} P_n z^n$.
\item Déterminer la fonction génératrice de $\displaystyle\sum_{n \ge 0} F_{n+2} F_n z^n$.
\item Montrer que
      $\displaystyle\sum_{n \ge 0} F_{n+1} F_n z^n = \frac{z}{1 - 2z - 2z^2 + z^3}$,
      sachant que
      $\displaystyle\sum_{n \ge 0} F_n^2 z^n = \frac{1 - z}{1 - 2z - 2z^2 + z^3}$.
\item Montrer que
      $\displaystyle\sum_{n \ge 0} F_{n+1} U_n(x) z^n = \frac{1 + 2xz}{1 - 2xz + (3 - 4x^2)z^2 + 4xz^3 + z^4}$,
      où $x = p_1 - p_2$, sachant que
      $\displaystyle\sum_{n \ge 0} F_n U_n(x) z^n = \frac{1 - xz}{1 - 2xz + (3 - 4x^2)z^2 + 4xz^3 + z^4}$.
\item Montrer que
      $\displaystyle\sum_{n \ge 0} F_n T_n(x) z^n = \frac{1 - xz + (1 - 2x^2)z^2}{1 - 2xz + (3 - 4x^2)z^2 + 2xz^3 + z^4}$.
\end{enumerate}
\textbf{Remarque :}
\begin{itemize}
\item la suite $(F_n)_{n \in \mathbb{N}}$ est donnée par $F_n = F_{n-1} + F_{n-2}$,
      $F_0 = 1$, $F_1 = 1$ ;
\item la suite $(P_n)_{n \in \mathbb{N}}$ est donnée par $P_n = 2P_{n-1} + P_{n-2}$,
      $P_0 = 0$, $P_1 = 1$ ;
\item la fonction génératrice d'une suite $(a_n)_{n \in \mathbb{N}}$ est
      $g(z) = \sum_{n \ge 0} a_n z^n$ ;
\item $T_n$ et $U_n$ sont respectivement les polynômes de Tchebychev de première et
      de seconde espèce.
\end{itemize}

\source{\emph{\small Université Mohamed Seddik Benyahia - Jijel, 2019, Épreuve~n\textsuperscript{o}~03 — barème : 6,5 points. Avertissement : le scan d'origine est de très mauvaise qualité ; l'identité introductive et les fractions des questions~3 à~5 doivent être recoupées avec le sujet papier.}}
\end{problem}

\newpage

\begin{problem}{(Jijel) — Épreuve N°03}{alg:jijel-2019-e3-p3}
Pour tous entiers naturels $n, r$ et $p$, les hyper-sommes de puissances d'entiers
notées $S_{p,(a,b)}^{(r)}(n)$ sont définies par la relation de récurrence
$$
S_{p,(a,b)}^{(r)}(n) = \sum_{k=0}^{n} S_{p,(a,b)}^{(r-1)}(k)
= \sum_{k=0}^{n} \binom{n + r - k}{r} (a + kb)^p, \quad n, r, p \ge 1,
$$
où
$S_{0,(a,b)}^{(r)}(n) = \sum_{k=0}^{n} \binom{n + r - k}{r} = \binom{n + r + 1}{r + 1}$
et $S_{p,(a,b)}^{(0)}(n) = \sum_{k=0}^{n} (a + kb)^p$. On note
$\binom{n}{k} = \dfrac{n!}{k!(n-k)!}$ le coefficient binomial.
\begin{enumerate}
\item Calculer la fonction génératrice exponentielle de $S_{p,(a,b)}^{(r)}(n)$ en $p$.
\item Si $r \equiv -1 \bmod n$, montrer que
      $S_{0,(a,b)}^{(r-1)}(n) \equiv (i + 1) \bmod n$ où $0 \le i \le n - 1$.
\item Montrer que pour tout nombre premier $n \ge 2$ et tout $p > 0$ on a
      $$
      S_{p,(a,b)}^{(r)}(n) \equiv (i + 2)\,a^p \bmod n
      \ \text{ si } \ r \equiv -1 \bmod n, \ \text{ où } \ 0 \le i \le n - 1.
      $$
\end{enumerate}

\source{\emph{\small Université Mohamed Seddik Benyahia - Jijel, 2019, Épreuve~n\textsuperscript{o}~03 — barème : 6,5 points. Avertissement : le texte d'origine est très pâle ; les conditions des questions~2 et~3 (l'indice $i$ et les membres de droite des congruences) doivent être recoupées avec le sujet papier.}}
\end{problem}
